\scanpage{169}
\mspage{383}

for the moment to ignore this possibility of normalization -- which brings out more clearly, indeed, the crucial character, in all these $l$-adic calculations, of the role played by the four quantities $c,d,t,u\in\widehat{\mathbb Z}$.

\scanpage{170}
\mspage{384}

\uline{45} \uline{The homomorphism} $\widetilde M_{0,3}\longrightarrow\mathrm{Gl}(2,\widehat{\mathbb Z})'$, \uline{construction of the extension} \deleted{\illegible{1}} $\widetilde N_{1,1}$ \uline{of} $\widetilde\Gamma$ \uline{by} $\mathrm{Sl}(2,\mathbb Z)^\wedge$.

We denote by $\widetilde M_{0,3}$ the group of \deleted{outer} loop automorphisms of $\unsure{\widehat{\mathcal G}_{0,3}^{+}}$ which \deleted{\unsure{commute}} \inserted{induce} up to conjugacy \deleted{with the action of $\unsure{\mathfrak S_3|\mathcal G_3}$, \unsure{which is}} \inserted{on the quotient $\unsure{\mathfrak S_3|\mathcal G_3}\simeq\widehat{\mathcal G}_{0,3}^{+}/\widehat\Pi_{0,3}$ the trivial action -- we have} thus a \inserted{structure of an} extension
\[
(1)\qquad 1\longrightarrow\unsure{\widehat{\mathcal G}_{0,3}^{+}}\longrightarrow\widetilde M_{0,3}\longrightarrow\widetilde N_{0,3}\longrightarrow 1,
\]
\margin{(We shall denote by $\widetilde M_{0,3}^{!}$ the subgroup consisting of the \unsure{automorphisms} which induce \unsure{strict\textsuperscript{ly}} \unsure{the ident\textsuperscript{ity}} in $\unsure{\mathfrak S_3|\mathcal G_3}$ -- it corresponds to $\widetilde N_{0,3}^{!}$ \unsure{-- automorphisms} that are outer loop automorphisms fixing the loop classes (and \unsure{\deleted{induce}|induce} the identity on $\unsure{\mathfrak S_3|\mathcal G_3}$ \unsure{up to conjug\textsuperscript{acy}}))}

$\widetilde N_{0,3}$ denotes the group of \uline{outer} \unsure{automorphisms} of $\widehat{\mathcal G}_{0,3}^{+}$ which \deleted{\illegible{3}} \inserted{induce the identity in} $\unsure{\mathfrak S_3|\mathcal G_3}$. We note that
\[
\Gamma=\Gamma_{\mathbb Q}\hookrightarrow\widetilde\Gamma,
\]
and that the inverse image of $\Gamma_{\mathbb Q}$ in the extension (1) is none other than $M_{0,3}$ -- but for the moment I have no very plausible characterization to propose of $\Gamma_{\mathbb Q}$ as a subgroup of $\widetilde\Gamma$. We have a canonical splitting of the extension (1), via the subgroup $\widetilde\Gamma_{\sigma}$ of $M$ consisting of the \unsure{automorphisms} $u\in\widetilde M_{0,3}$ which commute with $\sigma_\infty$,



\scanpage{171}

so that $\widetilde M_{0,3}$ presents itself as a $\frac12$ direct product
\[
(3)\qquad \widetilde M_{0,3}\simeq\underbrace{\widehat{\mathcal G}_{0,3}^{+}}_{\mathrm{Sl}(2,\mathbb Z)'{}^\wedge}\cdot\widetilde\Gamma_\sigma.
\]

Let us consider the canonical homomorphism
\[
(4)\qquad \unsure{\widehat{\mathcal G}_{0,3}^{+}}\simeq\mathrm{Sl}(2,\mathbb Z)'{}^\wedge\xrightarrow{\varphi_0'}\mathrm{Sl}(2,\widehat{\mathbb Z})'\subset\mathrm{Gl}(2,\widehat{\mathbb Z})'
\]
and the \deleted{homomorphism}\inserted{map}
\[
\begin{aligned}
(5)\qquad \widetilde\Gamma_\sigma&\xrightarrow{\varphi_\sigma'}\mathrm{Gl}(2,\widehat{\mathbb Z})'\\
\underline u&\longmapsto\mu(t\sigma_\infty+u)\tau_\infty^\varepsilon=\mu\begin{pmatrix}u&\varepsilon t\\-t&\varepsilon u\end{pmatrix},
\end{aligned}
\]
where \deleted{$\unsure{t,u,\alpha=\alpha_2}$ \illegible{2}}\deleted{\unsure{$\varepsilon=\varepsilon_2$}}
\[
t,u\in\widehat{\mathbb Z},\quad\inserted{\mu\in\widehat{\mathbb Z}^{*},}\quad\varepsilon=\varepsilon_2\in\{\pm1\},\quad\alpha=\alpha_2\in\mathbb Z/2\mathbb Z
\]
are the quantities (depending on $\underline u\in\widetilde\Gamma_\sigma$) considered in the \S\ above, with
\[
\varepsilon=(-1)^\alpha.
\]
If need be, we shall write $t(\underline u)$, $\underline u(u)$, $\unsure{\alpha(\underline u)}$, $\varepsilon(\underline u)$ to indicate the dependence \unsure{with respect} to $\underline u$.

\deleted{\uline{Theorem}}\inserted{Let us consider the (continuous) map}
\[
(6)\qquad\varphi':\widetilde M_{0,3}\longrightarrow\mathrm{Gl}(2,\widehat{\mathbb Z})'
\]
defined by
\[
(7)\qquad\varphi'(a.\underline u)=\varphi_0'(a)\varphi_\sigma'(\underline u)\qquad\left\{\begin{aligned}&a\in\widehat{\mathcal G}_{0,3}^{+}\\&\underline u\in\widetilde\Gamma_\sigma\end{aligned}\right.
\]

\scanpage{172}
\mspage{385}

\uline{Theorem} The map $\varphi'$ is a group homomorphism.

This means, taking into account that $\varphi_0$ is already a group homomorphism, the theorem means that

a) $\varphi_\sigma'$ is a group homomorphism, i.e.
\[
(8)\qquad\varphi_\sigma'(\underline u'\underline u)=\varphi_\sigma'(\underline u')\unsure{\,\varphi_\sigma'(\underline u)}\qquad(\underline u,\underline u'\in\widetilde\Gamma_\sigma)
\]
b) we have, for $a\in\unsure{\mathcal G_{0,3}^{+}|\widehat{\mathcal G}_{0,3}^{+}}$, $\underline u\in\widetilde\Gamma_\sigma$
\[
(9)\qquad\varphi_0'\bigl(\underbrace{\underline u a\underline u^{-1}}_{\underline u(a)}\bigr)=\varphi_\sigma'(\underline u)\varphi_0'(a)\varphi_\sigma'\bigl(\underline u^{\deleted{\unsure{-1}}}\bigr)^{-1}.
\]

\deleted{To verify the}\inserted{But the} formula (9) is a fact \inserted{already established,} \deleted{\unsure{consequence of (8)}, \illegible{1} the} since the \deleted{\illegible{1}} very construction of the matrix $\varphi_\sigma'(\underline u)=\mu(t\sigma_\infty+u)\tau_\infty^{\unsure{\varepsilon}}$ in the preceding paragraphs was made in a way \deleted{to make} \unsure{such} that this formula should be \unsure{verified} (cf. (67)(68), and the corollary to theorem (77)\ldots). To verify (8), \deleted{we must} \inserted{we shall} first make explicit the quantities $g,h,f$ for a product of two operators $\unsure{u',u''|u',u}\in\widetilde M_{0,3}$, in terms of their values for $u',u$ themselves. An immediate calculation gives, in this sense, some ``cocycle relations''
\[
(10)\qquad\left\{\begin{aligned}
g(\underline u'\underline u)&=\unsure{\underline u'}(g(\underline u)).g(\underline u')\\
h(\underline u'\underline u)&=\underline u'(h(\underline u)).h(\underline u')\\
f(\underline u'\underline u)&=\unsure{f(\underline u')}.\underline u'(f(\underline u))
\end{aligned}\right.
\]

\scanpage{173}

\inserted{\uline{Remarks} 1)} We shall moreover have likewise, with the notations of the \S\ above
\[
(11)\qquad\left\{\begin{aligned}
\alpha(u'u)&=u'(\alpha(u)).\alpha(u')\\
\beta(u'u)&=u'(\beta(u)).\beta(u')
\end{aligned}\right.
\]

\deleted{\unsure{(N.B.}} 2) The formulas (10) are a little more \unsure{friendly|friendly} for $f$, than for $g$ and $h$, \inserted{$\alpha$ and $\beta$} (-- one would feel like changing notation yet again, and also replacing $g$ and $h$ by $g^{-1},h^{-1}$ (as previously \unsure{a $g_0$} had become $f^{-1}$\ldots).

Let us again denote by the same letters $g(u),h(u)$ etc. the images of these quantities in $\mathrm{Sl}(2,\widehat{\mathbb Z})'$, \deleted{the formulas \unsure{(10)} \inserted{(11)} \unsure{give}} the analogous formulas in $\mathrm{Sl}(2,\widehat{\mathbb Z})'$ (or $(\mathrm{Gl}(2,\widehat{\mathbb Z})^{\pm1})'$ for \deleted{\unsure{(11)}} the $\alpha$ of (11)) which are written
\[
(12)\qquad\left\{\begin{aligned}
g(\underline u'\underline u)&=\mathrm{int}(\varphi_\sigma'(\underline u'))(g(\underline u)).g(\underline u')\\
h(\underline u'\underline u)&=\mathrm{int}(\varphi_\sigma'(\underline u'))(h(\underline u)).h(\underline u')\\
f(\underline u'\underline u)&=\deleted{\unsure{\mathrm{int}(\varphi_\sigma'(\underline u'))}}f(\underline u').\bigl(\mathrm{int}\unsure{(\varphi_\sigma(\underline u'))}(f(\underline u))\bigr)\\
\inserted{\alpha(\underline u'\underline u)}&\inserted{{}=\mathrm{int}\,\varphi_\sigma'(\underline u')(\alpha(\underline u)).\alpha(\underline u')}\\
\inserted{\beta(\underline u'\underline u)}&\inserted{{}=\mathrm{int}\,\varphi_\sigma'(\underline u')(\beta(\underline u)).\beta(\underline u')}
\end{aligned}\right.
\]
in $\mathrm{Gl}(2,\widehat{\mathbb Z})'$.

\deleted{To verify}\inserted{These formulas} in $\mathrm{Gl}(2,\widehat{\mathbb Z})'$ do indeed have an entirely explicit character from the point of view of effective computation, unlike

\scanpage{174}
\mspage{386}

the case of the initial formulas in $\widehat{\Pi}_{0,3}$ or $\unsure{\widehat{\widetilde{\Pi}}_{0,3}}$, where quantities like $\underline{u}'(\underbrace{g(\underline{u})}_{g})$ cannot be made explicit \deleted{\illegible{1}} \unsure{by} simple algebraic formulas in terms of the data \inserted{$\unsure{g,h,f,}$} $g',h',f'$ \deleted{\unsure{resp.}} $=g(\underline{u})$ etc. which determine $\underline{u}$ and $\underline{u}'$.

These formulas are supplemented, for the record, by
\[
\tag{13}
\mu(\underline{u}'\underline{u})=\mu(\underline{u}')\mu(\underline{u})
\]
implying in particular
\[
\tag{14}
\varepsilon_2(\underline{u}'\underline{u})=\varepsilon_2(\underline{u}')\varepsilon_2(\underline{u}).
\]

To carry out the calculation of $\varphi_{\sigma}(\underline{u}'\mathbin{.}\underline{u})$ and verify that it is equal to $\varphi_{\sigma}(\underline{u}')\varphi_{\sigma}(\underline{u})$, we \deleted{\unsure{shall have}} shall introduce the coefficients $r,s,t,u$ of $f=f(\underline{u})$, $r',s',t',u'$ of \inserted{$f'=$} $f(\underline{u}')$, and $r'',s'',t'',u''$ of $f''=f(\underline{u}'\underline{u})$, likewise \inserted{$\mu,\mu',\mu''$,} $\varepsilon,\varepsilon',\varepsilon''$, whence
\[
\tag{15}
\deleted{\text{and }\unsure{\text{relations}}}\qquad
\mu''=\mu\mu',\quad \varepsilon''=\varepsilon\varepsilon'
\]
and the formula in $f$ in (12) is therefore written
\[
\tag{16}
f''=f'\underline{u}'f\underline{u}'^{-1}
\]
where we have written, as in the \S\ above, $\underline{u}'$ instead of $\varphi_{\sigma}(\underline{u}')$, \deleted{\unsure{In}} In explicit terms
\[
\begin{pmatrix}r''&s''\\t''&u''\end{pmatrix}
=\begin{pmatrix}r'&s'\\t'&u'\end{pmatrix}
\unsure{\mu'}\begin{pmatrix}u'&\varepsilon't'\\-t'&\varepsilon'u'\end{pmatrix}
\begin{pmatrix}r&s\\t&u\end{pmatrix}
\unsure{\deleted{\frac{\mu'}{\mu'}}|\frac{\mu'}{\mu'}}
\begin{pmatrix}\varepsilon'u'&-\varepsilon't'\\t'&u'\end{pmatrix}
\]
\margin{N.B. $\det\underline{u}'=\mu'$}

\scanpage{175}

The product of the last three terms on the right-hand side is not encouraging:
\[
\underline{u}'f\underline{u}'^{-1}
=\mu'\deleted{\unsure{\frac{1}{\mu'}}}
\begin{pmatrix}
\varepsilon'ru'^2+st'u'+tt'u'+\varepsilon'ut'^2
&-\varepsilon'rt'u'+su'^2-tt'^2+\varepsilon'uu't'\\
-\varepsilon'ru't'\unsure{-}st'^2+tu'^2+\varepsilon'uu't'
&-\varepsilon'rt'^2\unsure{-}su't'-tu't'+\varepsilon'uu'^2
\end{pmatrix},
\]
but the product with $f'$ simplifies miraculously, thanks to the relation
\[
\unsure{t'^2}+u'^2=\varepsilon'/\mu'
\]
and we find
\[
\deleted{\unsure{(17)}}\qquad
f''=\begin{pmatrix}r''&s''\\t''&u''\end{pmatrix}
=\deleted{\unsure{\frac{1}{\mu'}}}
\begin{pmatrix}
\unsure{\varepsilon|*}\ \text{mess}&\text{mess}\\
\varepsilon'tu'+ut'&-\varepsilon'tt'+uu'
\end{pmatrix}
\]
whence
\[
\tag{18}
t''=\frac{1}{\mu'}(\varepsilon'tu'+ut')
\qquad
u''=\frac{1}{\mu'}(-tt'+uu')
\]
and hence (taking account of \inserted{$\mu''=\mu\mu'$,} $\varepsilon''=\varepsilon\varepsilon'$)
\[
\deleted{\unsure{(19)}}\qquad
\underline{u}''=\mu\mu'\deleted{\unsure{\mu'}}
\begin{pmatrix}
-\varepsilon'tt'+uu'&\varepsilon tu'+\varepsilon\varepsilon'ut'\\
-\varepsilon'tu'-ut'&-\varepsilon\varepsilon'tt'+\varepsilon\varepsilon'uu'
\end{pmatrix}
\]
and the formula to be checked \unsure{$u$},
\[
\underline{u}''=\underline{u}'\underline{u}
\]
reduces to
\[
\mu\mu'\begin{pmatrix}
-tt'+uu'&\varepsilon tu'+\varepsilon\varepsilon'ut'\\
-\varepsilon'tu'\deleted{\unsure{\varepsilon'}}-ut'&-\varepsilon\deleted{\unsure{\varepsilon'}}tt'+\varepsilon\varepsilon'uu'
\end{pmatrix}
=\mu'\begin{pmatrix}u'&\varepsilon't'\\-t'&\varepsilon'u'\end{pmatrix}
\mu\begin{pmatrix}u&\varepsilon t\\-t&\varepsilon u\end{pmatrix}
\]
which is indeed true, OK !.

Here again, one feels that these calculations ought to fit into a carpet of linear algebra that I am too lazy to unroll in extenso. One would have to paraphrase the group \inserted{\unsure{Or rather}} $\widetilde{M}_{0,3}$ -- described by these systems $(g,h,f)$ -- in terms \inserted{$(\alpha,\beta,f)$} of the reductions of $g,h,f$

\scanpage{176}
\mspage{387}

in $\mathrm{Gl}(2,\widehat{\mathbb Z})^{\pm1\,\prime}_{\unsure{\mathfrak g}}$, then replacing $\widehat{\mathbb Z}$ by an arbitrary ring, and of course $\mathrm{Gl}(2,\widehat{\mathbb Z})'$ by $\mathrm{Gl}(2,\unsure{A|\widehat{\mathbb Z}})$ itself. Perhaps I will come back to this eventually....

Having thus constructed
\[
\varphi:\widetilde{M}_{0,3}\longrightarrow\mathrm{Gl}(2,\widehat{\mathbb Z})'=\mathrm{Gl}(2,\widehat{\mathbb Z})/\mathord{\pm1}
\]
we deduce from it, by inverse image of the \inserted{central} extension $\mathrm{Gl}(2,\widehat{\mathbb Z})$ of $\mathrm{Gl}(2,\widehat{\mathbb Z})'$ by \deleted{$\unsure{\mathbb Z}$} $\{\pm1\}$, a central extension of $\widetilde{M}_{0,3}$ by $\{\pm1\}$, which I denote by $\unsure{\widetilde{N}_{1,1}}$, and we have a homomorphism of extensions
\[
\tag{19}
\diagram{\begin{tikzcd}[ampersand replacement=\&]
1\arrow[r]\&\{\pm1\}\arrow[r]\arrow[d,equal]\&\widetilde{N}_{1,1}\arrow[r]\arrow[d,"\varphi"]\&\deleted{\unsure{\mathcal T}}\widetilde{M}_{0,3}\arrow[r]\arrow[d,"\varphi'"]\&1\\
1\arrow[r]\&\{\pm1\}\arrow[r]\&\mathrm{Gl}(2,\widehat{\mathbb Z})\arrow[r]\&\mathrm{Gl}(2,\widehat{\mathbb Z})'\arrow[r]\&1
\end{tikzcd}}
\]

\deleted{On the other hand, denoting by \deleted{$\unsure{\widetilde{N}_{1,1}}$} $\unsure{\widetilde{\mathcal T}_{1,1}^{+}|\widehat{\mathcal T}_{1,1}^{+}}$ the inverse image of \deleted{$\unsure{\widetilde{N}_{1,1}}$} $\widehat{\mathcal G}_{0,3}^{+}\simeq\unsure{\widehat{\mathrm{Sl}(2,\mathbb Z)}'}\subset\widetilde{M}_{0,3}$ in the extension \deleted{$\unsure{\mathcal T}$} $\widetilde{N}_{1,1}$ of $\widetilde{M}_{0,3}$ by $\pm1$, we will have
\[
\diagram{\begin{tikzcd}[ampersand replacement=\&]
1\arrow[r]\&\unsure{\widehat{\mathcal T}_{1,1}^{+}|\widetilde{\mathcal T}_{1,1}^{+}}\arrow[r]\&\widetilde{N}_{1,1}
\end{tikzcd}}
\]
}

\scanpage{177}

One should check that the homomorphism $\varphi'$ we have just constructed coincides on $M_{0,3}\subset\widetilde{M}_{0,3}$ with the well-known ``modular'' homomorphism of \S\ 39 \unsure{(1)} -- it suffices to check it on the subgroup $\Gamma_{\mathbb Q,\sigma}$ of $\widetilde{\Gamma}_{\sigma}$ corresponding to $\Gamma_{\mathbb Q}$ -- this should not be very difficult, and I will \unsure{assume it} here. This \deleted{\unsure{being}} assumed, we find that the inverse image of $M_{0,3}$ in \deleted{$\unsure{\mathcal T}$} $\widetilde{N}_{1,1}$ is none other (up to canonical isomorphism) than $N_{1,1}$ itself -- in any case the inverse image of $\widehat{\mathcal G}_{0,3}^{+}$ \deleted{\unsure{on}} is none other than $\widehat{\mathcal T}_{1,1}^{+}$ \inserted{$\simeq\mathrm{Sl}(2,\mathbb Z)^\wedge$}. We will then have (with this identification)
\[
\tag{20}
\diagram{\begin{tikzcd}[ampersand replacement=\&]
1\arrow[r]\&\begin{gathered}\widehat{\mathcal T}_{1,1}^{+}\\\unsure{\Vert}\ \mathrm{Sl}(2,\mathbb Z)^\wedge\end{gathered}\arrow[r]\arrow[d]\&\deleted{\unsure{\mathcal G}}\widetilde{N}_{1,1}\arrow[r]\arrow[d,"\varphi"]\&\widetilde{\Gamma}\arrow[r]\arrow[d,"\unsure{\chi}"]\&1\\
1\arrow[r]\&\mathrm{Sl}(2,\widehat{\mathbb Z})\arrow[r]\&\mathrm{Gl}(2,\widehat{\mathbb Z})\arrow[r,"\det"]\&\widehat{\mathbb Z}^{*}\arrow[r]\&1
\end{tikzcd}}
\]

\scanpage{178}
\mspage{388}

so we have \unsure{reconstructed|reconstructed} in terms of loop \unsure{automorphisms} of $\widehat{\Pi}_{0,3}$ the diagram (1) of \S\ 39, or rather a diagram \unsure{a} larger, into which that of \S\ 39 embeds -- it would still remain to characterize $\Gamma_{\mathbb Q}$ as a subgroup of $\widetilde{\Gamma}$.

The diagram deduced directly from the preceding theorem
\[
\tag{21}
\diagram{\begin{tikzcd}[ampersand replacement=\&]
1\arrow[r]\&\widehat{\mathcal G}_{0,3}^{+}\arrow[r]\arrow[d]\&\widetilde{M}_{0,3}\arrow[r]\arrow[d,"\varphi'"]\&\widetilde{\Gamma}\arrow[r]\arrow[d,"\chi"]\&1\\
1\arrow[r]\&\mathrm{Sl}(2,\widehat{\mathbb Z})'\arrow[r]\&\unsure{\mathrm{Gl}(2,\widehat{\mathbb Z})'}\arrow[r,"\det"']\&\widehat{\mathbb Z}^{\unsure{*}}\arrow[r]\&1
\end{tikzcd}}
\]
can be regarded as being deduced from the preceding one by reduction modulo the \deleted{centers} \inserted{central subgroups} $\{\pm1\}$ in the two rows.

What becomes in (20) of the splitting (3) of $\widetilde{M}_{0,3}$? Is there a privileged way, for \unsure{every} element of $\widetilde{N}_{1,1}$ -- i.e. a pair of a $u=(g,h,f)\in\widetilde{M}_{0,3}$, and a $\underline{u}_0\in\mathrm{Gl}(2,\widehat{\mathbb Z})$

\scanpage{179}

which lifts the element $\underline{u}=\varphi'_{\unsure{\sigma}}(u)$ of $\mathrm{Gl}(2,\widehat{\mathbb Z})'$, to ``normalize'' it modulo \deleted{\illegible{1}} composition with an element of $\mathrm{Sl}(2,\mathbb Z)^\wedge$, so that it has \deleted{a} (say) a commutation property with $\sigma_\infty$, of the type $u(\sigma_\infty)=\deleted{\varepsilon}\sigma_\infty$, \deleted{\illegible{1}} $\unsure{u_0|u}(\sigma_\infty)=\pm\sigma_\infty$?

Another way of posing the question, in terms of the \unsure{subgroup} $\widetilde{\Gamma}_\sigma\subset\widetilde{M}_{0,3}$, is the following: is there a canonical lifting $\widetilde{\Gamma}_\sigma\longrightarrow\widetilde{N}_{1,1}$, thus assigning to every $\underline{u}=(g,f)\in\widetilde{\Gamma}_\sigma$ (\unsure{That is} a $\underline{u}=\mu(t\sigma_\infty+u)\tau_\infty^\alpha\in\mathrm{Gl}(2,\widehat{\mathbb Z})'$\unsure{,|\}} \deleted{\unsure{modulo}} i.e. a $\underline{u}$ defined modulo sign) a representative $\underline{u}_0\in\mathrm{Gl}(2,\widehat{\mathbb Z})$ -- i.e. a way of removing this ambiguity of sign? In fact, if we replace $\widetilde{\Gamma}_\sigma$ by $\Gamma_{\mathbb Q,\sigma}$, the Standard \deleted{\illegible{1}}\inserted{\unsure{yogas}} \unsure{suggest} that the liftings in

\scanpage{180}
\mspage{389}

question \unsure{should|must} correspond to elliptic curves over $\mathbb Q$, equipped with a Jacobi rigidification of level $2$, \unsure{corresponding} (under the isomorphism $\deleted{\unsure{\mathcal T}}\unsure{\mathcal M}_{1,1\,\mathbb Q}'\simeq U_{0,3\,\mathbb Q}$) to the point $1/2$ of $U_{0,3}(\mathbb Q)$. The indeterminacy in the choice of such an elliptic curve is precisely in $H^1(\mathbb Q,\unsure{\{\pm1\}})$ (it is a \unsure{datum} of a quadratic extension \unsure{of|to} $\mathbb Q$) which is identified with $\mathrm{Hom}(\Gamma_{\mathbb Q},\deleted{\unsure{\mu}}\{\pm1\})$ -- so this is indeed also the indeterminacy for a lifting. This proves \unsure{essentially} that the idea suggested by the general yoga is indeed \unsure{correct|correct} -- and this shows at the same time that there is no (even on $\Gamma_{\mathbb Q,\sigma}$) \unsure{really} privileged choice. That there are such liftings \inserted{(on $\Gamma_{\mathbb Q,\sigma}$)} comes from the fact that there are indeed elliptic

\scanpage{181}

curves answering the question, the simplest being given by the equation
\[
(22)\qquad y^{\unsure{2}}-x(x-1)(x-1/2)=0
\]
-- \unsure{and} there is the choice made ! It would remain to try to make it explicit \deleted{by} in a computational \unsure{formulary}, if this is not \unsure{hopeless}. I should examine this \deleted{question} question more closely, sometime.

\scanpage{182}
\mspage{390}

\uline{46}\quad \uline{Definitive notations} (?) -- \uline{new recapitulation}

(\uline{for} the loop automorphisms of $\widehat{\Pi}_{0,3}$ and of $\widehat{\mathcal G}_{0,3}^{+}\simeq\unsure{\mathrm{Sl}(2,\mathbb Z)^{\prime\wedge}}$).

\textcircled{I}\quad \uline{The groups} $\widetilde M_{0,3}$ \uline{and} $\widetilde N_{0,3}^{!}$.

We set

$(1)\qquad\widetilde M_{0,3}=$ group of automorphisms of $\widehat{\mathcal G}_{0,3}^{+}\simeq\unsure{\mathrm{Sl}(2,\mathbb Z)^{\wedge\prime}}$ which leave $\widehat{\Pi}_{0,3}$ invariant, and which respect the ``loop structure'' i.e. \unsure{satisfy} $u(\unsure{\varepsilon_i|\varepsilon_0})=\operatorname{int}(f^{-1})(\varepsilon_i^\mu)$, with $f\in\widehat{\mathcal G}_{0,3}^{+}$, $i\in\{0,1,\infty\}$, $\mu\in\widehat{\mathbb Z}^{*}$ (OBSERVATION $f\in\widehat{\Pi}_{0,3}$, and $\unsure{f|f_0}$ is then determined in $L_0\backslash\widehat{\Pi}_{0,3}$)

$\overset{\mathrm{restr.}}{\overset{\sim}{\longrightarrow}}$ group of \deleted{outer} loop automorphisms of $\widehat{\Pi}_{0,3}$ which commute as outer actions with the natural outer action of $\unsure{\mathfrak S_3|\mathcal G_3}$.

\margin{\uline{N.B.} $\widetilde M_{0,3}$ is also the normalizer of $\widehat{\mathcal G}_{0,3}^{+}$ in $\widehat{\widehat{\mathcal G}}_{0,3}$, and likewise $\widetilde N_{0,3}$ the \unsure{normalizer|normalizer} of $\unsure{\widehat{\mathcal T}_{0,3}}=\unsure{\mathcal T_{0,3}}$ in $\widehat{\widehat{\mathcal T}}_{0,3}$}

We have an extension structure
\[
(2)\qquad 1\longrightarrow\widehat{\Pi}_{0,3}\longrightarrow\widetilde M_{0,3}\longrightarrow\widetilde N_{0,3}\longrightarrow1,
\]
\unsure{where} \deleted{$\unsure{N_0}$}
\[
(3)\qquad\widetilde N_{0,3}\overset{\mathrm{dfn}}{=}\widetilde M_{0,3}/\widehat{\Pi}_{0,3}
\]
$\simeq$ group of \uline{outer} loop automorphisms of $\widehat{\Pi}_{0,3}$, commuting with $\unsure{\mathfrak S_3|\mathcal G_3}$.

The natural homomorphism: $\unsure{\mathfrak S_3|\mathcal G_3}\longrightarrow\operatorname{Aut}(\unsure{\mathfrak S_3|\mathcal G_3})$ is an \unsure{isomorphism}
\[
(4)\qquad\unsure{\mathfrak S_3|\mathcal G_3}\overset{x\mapsto\operatorname{int}(x)}{\overset{\sim}{\longrightarrow}}\operatorname{Aut}(\unsure{\mathfrak S_3|\mathcal G_3}),
\]
and the natural homomorphism
\[
\widetilde M_{0,3}\longrightarrow\operatorname{Aut}(\unsure{\mathfrak S_3|\mathcal G_3}=\widehat{\mathcal G}_{0,3}^{+}/\widehat{\Pi}_{0,3})
\]
then defines a homomorphism
\[
(5)\qquad\widetilde M_{0,3}\longrightarrow\unsure{\mathfrak S_3|\mathcal G_3}
\]
which can also be interpreted as the composite
\[
(6)\qquad\widetilde M_{0,3}\hookrightarrow
\underset{\substack{\text{all the loop automorphisms}\\\text{of }\widehat{\Pi}_{0,3}}}{\widehat{\widehat{\mathcal G}}_{0,3}}
\xrightarrow{\substack{\text{action on the}\\\text{loop classes}}}
\unsure{\mathfrak S_3|\mathcal G_3},
\]

\scanpage{183}

and which is trivial on $\widehat{\Pi}_{0,3}$, hence factors as
\[
(7)\qquad\widetilde N_{0,3}\longrightarrow\unsure{\mathfrak S_3|\mathcal G_3}.
\]
The kernels of (5), (7), are denoted $\widetilde M_{0,3}^{!}$, $\widetilde N_{0,3}^{!}$ respectively, so that we have a diagram of exact sequences

\unsure{\diagram{
\begin{tikzcd}[ampersand replacement=\&]
 \& \& \inserted{1}\arrow[d] \& \& \& \\
 \& \& \inserted{\Pi_{0,3}}\arrow[d] \& \& \& \\
(8) \& 1\arrow[r] \& \widetilde M_{0,3}^{!}\arrow[r]\arrow[d] \& \widetilde M_{0,3}\arrow[r]\arrow[d] \& \unsure{\mathfrak S_3|\mathcal G_3}\arrow[r]\arrow[d,equal] \& 1 \\
(9) \& 1\arrow[r] \& \widetilde N_{0,3}^{!}\arrow[r]\arrow[d] \& \widetilde N_{0,3}\arrow[r] \& \unsure{\mathfrak S_3|\mathcal G_3}\arrow[r] \& 1, \\
 \& \& \inserted{1} \& \& \&
\end{tikzcd}
}}

(9) being interpretable as the quotient of (8) by $\widehat{\Pi}_{0,3}$.

\inserted{\unsure{Let us note} also the \unsure{corresponding} extension structure
\[
(9\mathrm{bis})\qquad1\longrightarrow\widehat{\Pi}_{0,3}\longrightarrow\widetilde M_{0,3}^{!}\longrightarrow\widetilde N_{0,3}^{!}\longrightarrow1
\]}

We also have a canonical \unsure{homomorphism}
\[
(10)\qquad\unsure{\mathfrak S_3|\mathcal G_3}\simeq\widehat{\mathcal G}_{0,3}^{+}/\widehat{\Pi}_{0,3}\hookrightarrow\widetilde N_{0,3},
\]
whose composite with 7 is the identity, whence a canonical isomorphism
\[
(11)\qquad\widetilde N_{0,3}\simeq\unsure{\mathfrak S_3|\mathcal G_3}\times\widetilde N_{0,3}^{!}
\]
On the other hand, we also have an inclusion
\[
(12)\qquad
\underset{\substack{\unsure{\wr|\simeq}\\\unsure{\mathfrak S_3|\mathcal G_3}\times\mathbb Z/2\mathbb Z}}{\mathcal T_{0,3}}
\hookrightarrow\widetilde N_{0,3}
\]
which on the factor $\unsure{\mathfrak S_3|\mathcal G_3}\inserted{=\mathcal T_{0,3}^{+}}$ \unsure{coincides} with (10), and which \deleted{on} maps the nontrivial element $\tau$ of $\mathbb Z/2\mathbb Z\inserted{=\mathcal T_{0,3}^{!}}$ to the element, also denoted $\tau$, of the second factor $\widetilde N_{0,3}^{!}$ of (11).

\scanpage{184}
\mspage{391}

The composite homomorphism
\[
(13)\qquad \widetilde M_{0,3}\unsure{\hookrightarrow|\subset}
\unsure{\widehat{\widehat{\mathcal G}}_{0,3}}
\xrightarrow[\chi]{\text{\unsure{multiplier}}}
\widehat{\mathbb Z}^{*}
\]
\deleted{\unsure{gives|defines}} a homomorphism
\[
(14)\qquad \widetilde M_{0,3}\xrightarrow{\chi}\widehat{\mathbb Z}^{*}
\]
surjective, factoring through $\widetilde N_{0,3}$ and even \unsure{through} $\widetilde N_{0,3}/\unsure{\mathfrak S_3|\mathcal G_3}$

\diagram{\[
(15)\qquad
\begin{tikzcd}[ampersand replacement=\&]
\widetilde N_{0,3}\arrow[rr,"\chi"]\arrow[dr] \&\& \widehat{\mathbb Z}^{*} \\
\& \begin{gathered}\widetilde N_{0,3}/\unsure{\mathfrak S_3|\mathcal G_3}\\ \simeq\widetilde N_{0,3}^{!}\end{gathered}\arrow[ur] \&
\end{tikzcd}
,\]}

and giving rise to extension structures

\diagram{\[
(16)\qquad
\begin{tikzcd}[ampersand replacement=\&]
1\arrow[r] \& \widetilde M_{0,3}^{+}\arrow[r]\arrow[d] \& \widetilde M_{0,3}\arrow[r,"\chi"]\arrow[d] \& \widehat{\mathbb Z}^{*}\arrow[r]\arrow[d,equal] \& 1 \\
1\arrow[r] \& \widetilde N_{0,3}^{+}\arrow[r] \& \widetilde N_{0,3}\arrow[r,"\chi"] \& \widehat{\mathbb Z}^{*}\arrow[r] \& 1
\end{tikzcd}
,\]}

where we have set
\[
(17)\qquad
\left\{\begin{aligned}
\widetilde M_{0,3}^{+}&=\operatorname{Ker}(\widetilde M_{0,3}\to\widehat{\mathbb Z}^{*})\\
\widetilde N_{0,3}^{+}&=\operatorname{Ker}(\widetilde N_{0,3}\to\widehat{\mathbb Z}^{*})
\end{aligned}\right.
\]
-- the second row of (16) may be interpreted as the quotient of the first row by $\widehat\Pi_{0,3}$.

We have a canonical embedding
\[
(18)\qquad
\Gamma_{\mathbb Q}\overset{\mathrm{dfn}}{=}
\operatorname{Gal}\left(
\underbrace{\overline{\mathbb Q}_0}_{\inserted{\substack{\text{algebraic}\\\text{closure of }\mathbb Q\\\text{in }\mathbb C}}}
/\mathbb Q_0\right)
\unsure{\hookrightarrow|\subset}\widetilde N_{0,3}^{!}
\]
sending the \unsure{automorphism} of complex

\scanpage{185}

conjugation $\tau$ to the element already denoted $(\tau)$ (cf. (12)). The inverse image of \unsure{this|the} subgroup $\Gamma_{\mathbb Q}$ \deleted{in $\unsure{M}$}\inserted{of $\widetilde N_{0,3}^{!}$} under $\widetilde M_{0,3}\to\widetilde N_{0,3}^{!}\simeq\widetilde N_{0,3}/\unsure{\mathfrak S_3|\mathcal G_3}$, i.e. the inverse image of \deleted{$\unsure{\mathfrak S_3|\mathcal G_3}\times\Gamma_{\mathbb Q}\subset$}
\[
(19)\qquad
N_{0,3}\overset{\mathrm{dfn}}{=}
\unsure{\mathfrak S_3|\mathcal G_3}\times\Gamma_{\mathbb Q}
\subset\widetilde N_{0,3}
=\unsure{\mathfrak S_3|\mathcal G_3}\times\widetilde N_{0,3}^{!},
\]
is the extension \deleted{\unsure{of}} denoted $M_{0,3}$ elsewhere:

\margin{[We have a \deleted{\illegible{1}}\inserted{sub-}extension of \inserted{\unsure{the extension}} $\widetilde M_{0,3}$ of (2)]}

\diagram{\[
(20)\qquad
\begin{tikzcd}[ampersand replacement=\&]
1\arrow[r] \& \widehat\Pi_{0,3}\arrow[r]\arrow[d,"\unsure{\iota}"] \& M_{0,3}\arrow[r]\arrow[d,shift right=1.5pt]\arrow[d,no head,shift left=1.5pt,shorten <=7pt,shorten >=7pt] \& \overbrace{\unsure{\mathfrak S_3|\mathcal G_3}\times\Gamma_{\mathbb Q}}^{N_{0,3}}\arrow[r]\arrow[d,"{(\alpha,\beta)\to\beta}"] \& 1 \\
1\arrow[r] \& \widehat{\mathcal G}_{0,3}^{+}\arrow[r] \& M_{0,3}\arrow[r] \& \Gamma_{\mathbb Q}\arrow[r] \& 1
\end{tikzcd}
.\]}

Of course, it is in $M_{0,3}$ \inserted{and $N_{0,3}$} rather than in $\widetilde M_{0,3}$ \inserted{and $\widetilde N_{0,3}$} that we are interested -- hence in the remarkable subgroup $\Gamma_{\mathbb Q}$ of $\widetilde N_{0,3}^{!}$. We are going to ``determine'' by computation the group $\widetilde N_{0,3}^{!}$ -- and it will be a matter in the rest of the work of finding a maximum of necessary conditions on elements of $\widetilde N_{0,3}^{!}$ for them to come from $\Gamma_{\mathbb Q}$ -- \unsure{until} finding, if possible, a characterization (at least plausible, if not proved!) of the elements of $\Gamma_{\mathbb Q}$ -- i.e. of $\Gamma_{\mathbb Q}$ itself.

\scanpage{186}
\mspage{392}

\Needspace{8\baselineskip}
II \uline{Computational determination of the elements of} $\widetilde N_{0,3}^{!}$.\par\nobreak
1) \uline{Determination by} $h,g,f$.

\deleted{The \unsure{specification} of $u\in\unsure{\widetilde M_{0,3}^{!}|\widetilde N_{0,3}^{!}}$}

Recall that
\[
\begin{aligned}
(21)\qquad
&\unsure{\widetilde N_{0,3}^{!}|\widetilde M_{0,3}^{!}}
=\text{automorphisms of }\widehat{\mathcal G}_{0,3}^{+}
\simeq\unsure{\mathrm{Sl}(2,\mathbb Z)^{\wedge\prime}|\widehat{\mathrm{Sl}(2,\mathbb Z)^{\prime}}}
\text{ that stabilize}\\
&\widehat\Pi_{0,3},\quad\text{induce the identity on }
\unsure{\mathfrak S_3|\mathcal G_3}
=\widehat{\mathcal G}_{0,3}^{+}/\widehat\Pi_{0,3},\\
&\text{and preserve the loop structure (defined}\\
&\text{by the class of }\varepsilon_0\text{)}
\end{aligned}
\]

A $u\in\widetilde M_{0,3}^{!}$ is determined by a system
\[
(22)\qquad (h,g,f)\quad
\deleted{\unsure{\in(\widehat{\Pi}_{0,3})^{3}}\,\unsure{\times}\,\unsure{\widehat{\Pi}_{0,3}}}
\qquad(h,g,f\in\widehat\Pi_{0,3})
\]
\margin{$f$ determined modulo $f\mapsto\ell_0^{n}f$, $n\in\widehat{\mathbb Z}$}
\deleted{\unsure{satisfying}}\inserted{via} the formulas
\[
(23)\qquad
\left\{\begin{aligned}
u(\sigma_\infty)&=h\sigma_\infty\\
u(\rho)&=g\rho
\end{aligned}\right.
\]
\deleted{$\unsure{u(\varepsilon_0)}=\unsure{\operatorname{int}}\unsure{(}$}

the condition that these formulas define an automorphism (or at least an endomorphism) of $\widehat{\mathcal G}_{0,3}^{+}$ (which will necessarily induce the identity on $\unsure{\mathfrak S_3|\mathcal G_3}$) is expressed by the formulas
\[
(24)\qquad (h\sigma_\infty)^{2}=(g\rho)^{3}=1,
\]
which can also be written
\[
(25)\qquad
\boxed{\left\{\begin{aligned}
h\sigma_\infty(h)&=1\\
g\rho(g)\rho^{\unsure{2}}(g)&=1
\end{aligned}\right.},
\]
and the condition that $u$ preserves the loop structure is expressed by
\[
(26)\qquad u(\varepsilon_0)=\operatorname{int}(f^{-1})(\varepsilon_0^{\mu})
\qquad \mu\in\widehat{\mathbb Z}^{*}\text{ suitable},
\]
\margin{loop relation in $\widehat{\mathcal G}_{0,3}^{+}$ (raw form)}
\unsure{equivalent|equivalent} to the condition

\scanpage{187}

\margin{loop relation in $\widehat\Pi_{0,3}$ (raw form)}
\[
(27)\qquad
u(\ell_0)=\operatorname{int}(f^{-1})(\ell_0^{\mu}),
\]
it is easily made explicit in terms of $h,g,f$ (preferably in form (26)) taking into account that $\varepsilon_0=\sigma_\infty\rho$
\[
\underbrace{(h\sigma_\infty)(g\rho)}_{h\sigma_\infty(g)\underbrace{\sigma_\infty\rho}_{\varepsilon_0}}
=\underbrace{\operatorname{int}(f^{-1})\varepsilon_0^{\mu}}_{f^{-1}\varepsilon_0^{\mu}f},\quad\text{i.e}
\]
\[
f h\sigma_\infty(g)\underbrace{\varepsilon_0 f^{-1}}_{\unsure{\varepsilon_0(f^{-1})}\varepsilon_0=\sigma_\infty(\rho(f)^{-1})\varepsilon_0}
=\varepsilon_0^{\mu}\qquad\text{i.e.}
\]
\[
f h\sigma_\infty(g\rho(f)^{-1})=\varepsilon_0^{\mu-1}
\qquad\text{i.e., setting}
\]
\[
(28)\qquad \unsure{\nu}\overset{\mathrm{dfn}}{=}\frac{\mu-1}{2}\in\widehat{\mathbb Z}
\]
\margin{possible, since $\mu\in\widehat{\mathbb Z}^{*}$ hence $\mu\equiv1\ (2)$}
the relation (taking into account that $\varepsilon_0^{2}=\ell_0$)
\[
f h\sigma_\infty(g\rho(f)^{-1})=\ell_0^{\nu},
\]
or again, applying $\sigma_\infty$, since $\sigma_\infty(\ell_0)=\ell_1$
\[
\sigma_\infty(f)\sigma_\infty(h)g\rho(f)^{-1}=\ell_1^{\nu}
\]
that is, finally
\[
(29)\qquad
\boxed{\sigma_\infty(h)g=\sigma_\infty(f)^{-1}\ell_1^{\nu}\rho(f)}
\]
(\uline{loop relation in} $h,g,f$).

\margin{\uline{N.B.} $\nu$\inserted{$\in\widehat{\mathbb Z}$} is uniquely determined by $g,h$, and \unsure{we have} also $2\nu+1\in\widehat{\mathbb Z}^{*}$ $2\nu+1$ is the multiplier\ldots}

\uline{N.B.} One must in addition require, for $h,g,f$ to really define an \uline{automorphism}, that (24) be a \uline{presentation-generation} of $\widehat{\mathcal G}_{0,3}^{+}$ -- in what follows, we shall always tacitly take this ``preliminary condition'' as understood.
